\documentclass[12pt]{article}
\usepackage[utf8]{inputenc}
\usepackage[T1]{fontenc}
\usepackage{hyperref}
\title{Semnam Language Resources for K-12 Curriculum Development}
\author{Compiled by Claude Sonnet 5.5 (claude-sonnet-5-5) for Andr\'as Kornai}
\date{}
\begin{document}
\maketitle

\section{Introduction}
Semnam (ISO 639-3 ssm) is an Austroasiatic (Aslian) language of Peninsular Malaysia, spoken in the Perak valley by a Lanoh subgroup near Gerik and Lenggong. The Lanoh are officially classified as Orang Asli. Speaker estimates differ: Burenhult and Wegener report some 250 speakers, while Wikipedia cites 670. The language was first described in 2009.

\section{Primary Educational Resources}
No primers, textbooks or school materials in Semnam turned up in the search.

\section{Linguistic and Reference Materials}
\begin{itemize}
\item Burenhult and Wegener (2009), Preliminary notes on the phonology, orthography and vocabulary of Semnam, Journal of the Southeast Asian Linguistics Society 1:283--312. It includes a proposed practical orthography and a wordlist of about 1,250 items: \url{https://www.mpi.nl/publications/item68372/preliminary-notes-phonology-orthography-and-vocabulary-semnam-austroasiatic}
\item Other records of the same paper: \url{https://pub.uni-bielefeld.de/record/2900407} and \url{https://lucris.lub.lu.se/ws/files/6330526/1790595.pdf}
\item Wikipedia, Semnam language: \url{https://en.wikipedia.org/wiki/Semnam_language}
\end{itemize}

\section{Religious and Community Resources}
No religious or community resources turned up in the search.

\section{Digital and Audio Resources}
No audio, video or software resources turned up in the search.

\section{Curriculum Development Context}
The only substantial resource found is the 2009 data paper, with its orthography proposal and wordlist. These could be a starting point for community-led materials, but no teaching materials exist in the search results. Any curriculum would need to be built with the community and with the Malaysian Orang Asli education context in mind.
\end{document}